\documentclass{article}
\usepackage{decomp}
\begin{document}
\sourcefile{poc4045-sup-0001-supporting\_information.docx}
\decompnote{The Supporting Information is a Word document (Wiley file name poc4045-sup-0001-supporting\_information.docx), not a PDF, so there are no page numbers to mark; the document has three hard page breaks (after the affiliations, after Fig. S3 and after the derivation), which are not marked. Text converted with pandoc and checked against the .docx; the equations are MathType 6 objects (OLE ``Equation Native'', MTEF v5) and were transcribed from their MTEF data (read with work/mtef\_dump.py) and checked against the WMF previews. The figures are the three embedded TIFF images (figures/sfig1.png, sfig2.png, sfig3.png, converted losslessly).}

\noindent Supporting Information for\par

\section*{Kinetic Detection of Osmium(VI) Ester Intermediates during the OsO\textsubscript{4}–Mediated Aqueous Dihydroxylation of Chloroethylenes}

\noindent Mario Gabričević\textsuperscript{a}, Gábor Lente\textsuperscript{b*} and István Fábián\textsuperscript{c,d}\par

\begin{quote}\small
\noindent \textsuperscript{a}Faculty of Pharmacy and Biochemistry, University of Zagreb, A. Kovačića 1, 10000 Zagreb, Croatia.\par
\noindent \textsuperscript{b}Department of General and Physical Chemistry, University of Pécs, Pécs, Hungary\par
\noindent \textsuperscript{c}Department of Inorganic and Analytical Chemistry, University of Debrecen, Debrecen, Hungary\par
\noindent \textsuperscript{d}MTA-DE Redox and Homogeneous Catalytic Reaction Mechanisms Research Group, Debrecen, Hungary\par
\end{quote}

\begin{center}
\includegraphics[width=0.75\linewidth,keepaspectratio]{figures/sfig1.png}
\end{center}
\noindent\textbf{Fig. S1.} Spectral change of OsO\textsubscript{4} with \emph{t}-DCE reaction. t = 10 s (top spectrum), t = 5 h (bottom spectrum). Conditions: [Os(VIII)] = 0.75 mmol/dm\textsuperscript{3}, [\emph{t}-DCE] = 8 mmol/dm\textsuperscript{3}, [NaOH] = 0.1 mol/dm\textsuperscript{3}, \emph{I} = 0.1 mol/dm\textsuperscript{3}, \emph{T} = 25 °C.
\alttext{Line plot of two absorption spectra. x axis ``$\lambda$ (nm)'', linear, 260 to 550 (left frame unlabelled), labelled 300--550 in steps of 50; y axis ``\emph{A}'', linear, 0.0 to about 1.8, labelled 0.0, 0.5, 1.0, 1.5. No legend; the caption says the top spectrum is t = 10 s and the bottom one t = 5 h. Red (top): A ≈ 1.65 at 260 nm, minimum ≈ 1.28 near 281 nm, broad maximum ≈ 1.52 near 313 nm, then falling (≈ 1.13 at 345 nm, ≈ 0.51 at 400 nm, ≈ 0.22 at 450 nm, ≈ 0.10 at 500 nm) to ≈ 0.06 at 550 nm. Blue (bottom): ≈ 1.6 at 260 nm, minimum ≈ 0.89 near 279 nm, maximum ≈ 1.07 near 299 nm, a shoulder at ≈ 0.5 around 335--350 nm, ≈ 0.12 at 400 nm, flat at ≈ 0.07 from 450 nm; it meets the red line at the right edge (≈ 0.06 at 550 nm). A thin black arrow points from the red maximum (≈ 315 nm) down to the blue maximum (≈ 300 nm).}
\figuredata{(single panel)}{figures/sfig1.csv}{figures/sfig1.py}{digitized (tier 3) from the embedded lossless TIFF (893 x 631 px, 220 ppi) by figures/sfig1\_extract.py}

\begin{center}
\includegraphics[width=0.75\linewidth,keepaspectratio]{figures/sfig2.png}
\end{center}
\noindent\textbf{Fig. S2.} Stoichiometry of the redox reaction between OsO\textsubscript{4} and \emph{c}-DCE. All spectra were measured after completion of the reaction (24 h). Conditions: [Os(VIII)] = 1 mmol/dm\textsuperscript{3}, [cDCE] = 0.05 – 5 mmol/dm\textsuperscript{3} [NaOH] = 0.10 mol/dm\textsuperscript{3}, \emph{I} = 0.10 mol/dm\textsuperscript{3}, \emph{T} = 25 °C.
\alttext{Line plot of absorption spectra with an inset. Main panel: x axis ``$\lambda$ (nm)'', linear, 260 to 500 (left frame unlabelled), labelled 300--500 in steps of 50; y axis ``\emph{A}'', linear, 0.0 to 2.0, labelled 0.0--2.0 in steps of 0.5. Ten spectra in different colours, no legend. Black (top, well apart): ≈ 1.68 at 262 nm, minimum ≈ 1.33 near 282 nm, maximum ≈ 1.72 near 319 nm with a small notch near 326 nm, then falling to ≈ 0.62 at 400 nm and ≈ 0.10 at 500 nm. The other nine (blue, red, dark cyan, magenta, wine, olive, pink, green, navy) rise steeply below about 270 nm, have a minimum ≈ 0.89--1.03 near 279--280 nm and a maximum ≈ 1.09--1.21 near 300 nm; between 320 and 400 nm they separate a little: blue and red highest (≈ 0.75 and 0.70 at 345 nm, with a shoulder), dark cyan (≈ 0.58), and a bundle of six at ≈ 0.48--0.54; at 500 nm all nine lie between ≈ 0.03 and 0.06. Inset (upper right, own frame): x axis ``[\emph{cis}-DCE]/[Os(VIII)]'', 0 to about 5.5, labelled 0--5; y axis ``A (345 nm)'', 0.4 to 1.5, labelled 0.4, 0.7, 1.0, 1.3. Nine grey diamonds: ≈ 1.45 at ratio ≈ 0.1, ≈ 0.76 and 0.72 at ≈ 0.3 and 0.5 (overlapping), ≈ 0.59 at 0.75, ≈ 0.54 at 1, ≈ 0.52 at 1.5, ≈ 0.50 at 2 and 4, ≈ 0.49 at 5. A red dotted straight line falls from ≈ 1.39 near ratio 0.06 to the bottom of the inset near ratio 1.15; an orange dashed line runs almost horizontally from ≈ 0.52 at ratio 0.5 to ≈ 0.49 at ratio 4.8.}
\figuredata{main panel}{figures/sfig2main.csv}{figures/sfig2main.py}{digitized (tier 3) from the embedded TIFF by figures/sfig2\_extract.py; 10 spectra named by line colour (no legend), values in the bundle flagged `overlap'}
\figuredata{inset}{figures/sfig2inset.csv}{figures/sfig2inset.py}{digitized (tier 3) from the same image by figures/sfig2\_extract.py: 9 markers and the two straight lines (end points of a least-squares line through their pixels)}

\begin{center}
\includegraphics[width=0.75\linewidth,keepaspectratio]{figures/sfig3.png}
\end{center}
\noindent\textbf{Fig. S3.} Stoichiometry of the redox reaction between OsO\textsubscript{4} and t-DCE. All spectra were measured after completion of the reaction (24 h). Conditions: [Os(VIII)] = 1 mM, [\emph{t}-DCE] = 0.05 – 5 mmol/dm\textsuperscript{3} [NaOH] = 0.10 mol/dm\textsuperscript{3}, I = 0.1 mol/dm\textsuperscript{3}, \emph{T} = 25 °C.
\alttext{Line plot of absorption spectra with an inset. Main panel: x axis ``$\lambda$ (nm)'', linear, 250 to 500, labelled 250--500 in steps of 50; y axis ``\emph{A}'', linear, 0.0 to about 1.85, labelled 0.0, 0.5, 1.0, 1.5. Eleven spectra in different colours, no legend. All cross close together near 262 nm at A ≈ 0.85--0.9; to the left they rise to between ≈ 1.1 (royal blue) and above the top frame. From top to bottom near 320 nm: royal blue (maximum ≈ 0.86 near 315 nm), green (≈ 0.81 near 308 nm), pink (≈ 0.71 near 304 nm), wine (≈ 0.69 near 300 nm), navy (≈ 0.66), olive (≈ 0.64), and a bundle of magenta, dark cyan, blue, red and black (maxima ≈ 0.61--0.65 near 300 nm, minima ≈ 0.51--0.55 near 280 nm, a shoulder ≈ 0.3 around 340--350 nm, the black line lowest). All decrease to ≈ 0.05--0.09 at 500 nm. Inset (upper right, own frame): x axis ``[\emph{trans}-DCE]/[Os(VIII)]'', 0 to about 5.4, labelled 0--5; y axis ``A (nm)'' (no wavelength given), about 0.27 to 0.87, labelled 0.4, 0.6, 0.8. Eleven grey diamonds: ≈ 0.82 at ratio ≈ 0.1, 0.75 at 0.2, 0.59 at 0.3, 0.51 at 0.4, 0.47 at 0.5 (the last two overlapping), 0.41 at 0.7, 0.37 at 1, 0.35 at 2, 0.33 at 3, 0.32 at 4, 0.31 at 5. A red dotted straight line falls from ≈ 0.80 near ratio 0.08 to the bottom of the inset near ratio 0.8; an orange dashed line falls gently from ≈ 0.37 at ratio 0.5 to ≈ 0.31 at ratio 4.8.}
\figuredata{main panel}{figures/sfig3main.csv}{figures/sfig3main.py}{digitized (tier 3) from the embedded TIFF by figures/sfig3\_extract.py; 11 spectra named by line colour (no legend), values in the bundle flagged `overlap'}
\figuredata{inset}{figures/sfig3inset.csv}{figures/sfig3inset.py}{digitized (tier 3) from the same image by figures/sfig3\_extract.py: 11 markers and the two straight lines}

\subsection*{Derivation of Eq. 14}
\decompnote{The heading says Eq. 14, but the section derives Eqs. 12 and 13 of the main article (Eq. 14 is the expression for \emph{k}\textsubscript{obs} of the fast process).}

The rate equation to be solved is given in Eq. 11 of the main manuscript:

\begin{equation}
\frac{d[\mathrm{EO}]}{dt} = k_+\bigl([\mathrm{E}]_0 - [\mathrm{EO}]\bigr)\bigl([\mathrm{O}]_0 - [\mathrm{EO}]\bigr) - k_-[\mathrm{EO}] \tag{11}
\end{equation}

The technique that can provide the solution of this ordinary differential equation can be found in the literature [37]. Here, only the final solution will be given:
\decompnote{The SI has no reference list of its own. In the main article's list, [37] is Y. E. Yan, F. W. Schwartz, Environ. Sci. Technol. 2000 (on chloroethylene oxidation); the main text cites [45] (G. Lente, Deterministic Kinetics in Chemistry and Systems Biology) for the solution of this rate equation.}

\begin{equation}
\begin{gathered}
[\mathrm{EO}] = \frac{A}{2k_+} - \frac{B\bigl(A + B + Ae^{-Bt} - Be^{-Bt}\bigr)}{2k_+\bigl(A + B - Ae^{-Bt} + Be^{-Bt}\bigr)}\\[4pt]
A = k_+[\mathrm{E}]_0 + k_+[\mathrm{O}]_0 + k_- \qquad B = \sqrt{A^2 - 4{k_+}^2[\mathrm{E}]_0[\mathrm{O}]_0}
\end{gathered} \tag{12}
\end{equation}

However, this solution and its resemblance to an exponential curve is a pivotal point in the kinetic analysis in this manuscript. Therefore, it will be proved that Eq. 12 is in fact a solution of Eq. 11 (for this derivation, the method for finding the formula of the solution is irrelevant). So the first derivative of the function shown in Eq. 12 will be calculated with respect to time:

\begin{equation}
\frac{d[\mathrm{EO}]}{dt} = \frac{B^2(A - B)e^{-Bt}}{2k_+\bigl(A + B - Ae^{-Bt} + Be^{-Bt}\bigr)} + \frac{B^2\bigl(A + B + Ae^{-Bt} - Be^{-Bt}\bigr)}{2k_+\bigl(A + B - Ae^{-Bt} + Be^{-Bt}\bigr)^2}(A - B)e^{-Bt} \tag{S1}
\end{equation}

Also, the solution shown in Eq. 12 can be re-arranged to the following form

\begin{equation}
\begin{gathered}
\frac{B}{A - 2k_+[\mathrm{EO}]} = \frac{A + B - Ae^{-Bt} + Be^{-Bt}}{A + B + Ae^{-Bt} - Be^{-Bt}} = \frac{A + B + Ae^{-Bt} - Be^{-Bt} - 2Ae^{-Bt} + 2Be^{-Bt}}{A + B + Ae^{-Bt} - Be^{-Bt}} =\\[4pt]
1 - 2\frac{(A - B)e^{-Bt}}{A + B + (A - B)e^{-Bt}}
\end{gathered} \tag{S2}
\end{equation}

Rearranging Eq. S2 gives a useful formula the use of which shortens the derivation:

\begin{equation}
\frac{(A - B)e^{-Bt}}{A + B + (A - B)e^{-Bt}} = \frac{1}{2} - \frac{B}{2\bigl(A - 2k_+[\mathrm{EO}]\bigr)} \tag{S3}
\end{equation}

Substituting this formula back into Eq. S2 and further transformations give:

\begin{equation}
\begin{split}
\frac{d[\mathrm{EO}]}{dt} &= \left(\frac{A}{2k_+} - [\mathrm{EO}]\right)\frac{B(A - B)e^{-Bt}}{A + B + (A - B)e^{-Bt}} + \left(\frac{A}{2k_+} - [\mathrm{EO}]\right)^2\frac{2k_+}{A + B + (A - B)e^{-Bt}}(A - B)e^{-Bt} =\\[4pt]
&\left(\frac{A}{2k_+} - [\mathrm{EO}]\right)\frac{(A - B)e^{-Bt}}{A + B + (A - B)e^{-Bt}}(B + A - 2k_+x) =\\[4pt]
&\frac{1}{2k_+}\bigl(A - 2k_+[\mathrm{EO}]\bigr)\left(\frac{1}{2} - \frac{B}{2\bigl(A - 2k_+[\mathrm{EO}]\bigr)}\right)\bigl(B + A - 2k_+[\mathrm{EO}]\bigr) =\\[4pt]
&\frac{1}{4k_+}\bigl(A - B - 2k_+[\mathrm{EO}]\bigr)\bigl(B + A - 2k_+[\mathrm{EO}]\bigr) =\\[4pt]
&\frac{1}{4k_+}\bigl(A^2 - B^2\bigr) + \frac{1}{2}(-A + B - B - A)[\mathrm{EO}] + k_+[\mathrm{EO}]^2 =\\[4pt]
&\frac{1}{4k_+}\bigl(A^2 - B^2\bigr) - A[\mathrm{EO}] + k_+[\mathrm{EO}]^2
\end{split} \tag{S4}
\end{equation}
\decompnote{In the second line of Eq. S4 the last factor reads (\emph{B} + \emph{A} − 2\emph{k}\textsubscript{+}\emph{x}); the next line has [EO] in place of \emph{x}. Transcribed as printed (MTEF data and WMF preview agree).}

From Eq. 12, the following equation can be deduced:

\begin{equation}
A^2 - B^2 = A^2 - \left(\sqrt{A^2 - 4{k_+}^2[\mathrm{E}]_0[\mathrm{O}]_0}\right)^2 = 4{k_+}^2[\mathrm{E}]_0[\mathrm{O}]_0 \tag{S5}
\end{equation}

Therefore, the last line of Eq. S4 is identical to the right hand side of Eq. 11, which proves that Eq. 12 is in fact the solution of Eq. 11.

For proving that Eq. 13 is a good approximation under conditions when the equilibrium is heavily shifted toward the initial substances, the exact solution is first transformed into an equivalent form:

\begin{equation}
\begin{split}
[\mathrm{EO}] &= \frac{A}{2k_+} - \frac{B\bigl(A + B + Ae^{-Bt} - Be^{-Bt}\bigr)}{2k_+\bigl(A + B - Ae^{-Bt} + Be^{-Bt}\bigr)} = \frac{A}{2k_+} - \frac{B\bigl(A + B - Ae^{-Bt} + Be^{-Bt} + 2Ae^{-Bt} - 2Be^{-Bt}\bigr)}{2k_+\bigl(A + B - Ae^{-Bt} + Be^{-Bt}\bigr)} =\\[4pt]
&= \frac{A}{2k_+} - \frac{B}{2k_+} - \frac{B(A - B)e^{-Bt}}{k_+\bigl(A + B - Ae^{-Bt} + Be^{-Bt}\bigr)}
\end{split} \tag{S6}
\end{equation}

As pointed out in the main text, the inequality \emph{k}\textsubscript{−} >> \emph{k}\textsubscript{+}[E]\textsubscript{0} + \emph{k}\textsubscript{+}[O]\textsubscript{0} holds in this system (the chemical meaning of this inequality is that the conversion is quite low in equilibrium, which is shifted toward the initial substances. Given the definitions of A and B in Eq. 12, the following inequality is also true:

\begin{equation}
A + B \gg Ae^{-Bt} - Be^{-Bt} \tag{S7}
\end{equation}

Therefore, the exponential terms in the denominator of Eq. S6 can be dropped:

\begin{equation}
[\mathrm{EO}] \cong \frac{A}{2k_+} - \frac{B}{2k_+} - \frac{B(A - B)e^{-Bt}}{k_+(A + B)} \tag{S8}
\end{equation}

Furthermore, \emph{k}\textsubscript{−} >> \emph{k}\textsubscript{+}[E]\textsubscript{0} + \emph{k}\textsubscript{+}[O]\textsubscript{0} also guarantees that A $\cong$ B $\cong$ \emph{k}\textsubscript{−}. With these, further simplifications can be made in Eq. S8:

\begin{equation*}
[\mathrm{EO}] \cong \frac{A}{2k_+} - \frac{B}{2k_+} - \frac{B(A - B)e^{-Bt}}{k_+(A + B)} \cong \frac{A}{2k_+} - \frac{B}{2k_+} - \frac{(A - B)e^{-Bt}}{2k_+} \cong \frac{(A - B)\bigl(1 - e^{-k_-t}\bigr)}{2k_+}
\end{equation*}
\decompnote{This last equation has no number in the SI.}

If the values of parameters A and B are substituted back into Eq. S8, the form given in Eq. 13 in the main article is obtained.

\subsection*{Detailed thermodynamic analysis of the redox reaction between osmium tetroxide and ethylene and its consequences for the Gibbs energy of formation of the osmate ester}

Although the experimental study was carried out for chloroethylenes, this section shows a detailed thermodynamic analysis of the reaction between ethylene and osmium tetroxide for three reasons:

\begin{hangpars}
1. All thermodynamic data are experimentally available for ethylene and its oxidized products

2. Quantum chemical calculations have been published for unsubstituted ethylene

3. The analysis is not influenced by the removal of chlorine atoms, which occurs after dihydroxylation for chloroethyelenes
\end{hangpars}

The overall reaction of the dihydroxylation of ethylene in its reaction with osmium tetroxide has the following stoichiometry:

\begin{center}
OsO\textsubscript{4} + C\textsubscript{2}H\textsubscript{4} + 2H\textsubscript{2}O → OsO\textsubscript{4}\textsuperscript{2−} + C\textsubscript{2}H\textsubscript{6}O\textsubscript{2} + 2H\textsuperscript{+}
\end{center}

As the experimental study was carried out in water as a solvent, relevant aqueous thermodynamic data will be used as a starting point of the analysis.

The following aqueous standard two-electron reduction potential is known experimentally:

\begin{center}
OsO\textsubscript{4} + 2 H\textsubscript{2}O + 2 e\textsuperscript{−} → OsO\textsubscript{2}(OH)\textsubscript{4}\textsuperscript{2−} \qquad \emph{E}° = 0.46 V
\end{center}

\noindent Sources:

D. Richens, \emph{The Chemistry of Aqua Ions}, Wiley, Chichester \textbf{1997}, p. 423.

R.N. Goldberg, L. G. Hepler, \emph{Chem. Rev.} \textbf{1968}, \emph{68}, 229.

W. M. Latimer, \emph{Oxidation States of the Elements and their Potential in Aqueous Solution}, Prentice-Hall, Englewood Clipps, \textbf{1952}, p. 202.

At first sight, this reduction potential may seem very low and it might also be somewhat counterintuitive that OsO\textsubscript{4} is only a mild oxidant in water. As a qualitative confirmation of this weak oxidizing power, we showed experimentally that aqueous OsO\textsubscript{4} does not oxidize iodide ion (\emph{E}° = 0.54 V) at all. Both OsO\textsubscript{4} and iodide ion are known to be highly labile in redox reactions, the reason for the lack of reaction can only be thermodynamic here, which clearly confirms the fact that OsO\textsubscript{4} is a mild oxidant.

All other thermodynamic data at were taken from the CRC Handbook of Chemistry and Physics (Ed. D. R Lide), 87th Edition, 1996 and are quoted for 298.15 K.

First, the standard Gibbs free energy of formation for gaseous ethylene is known: C\textsubscript{2}H\textsubscript{4} (g) ∆\textsubscript{f}\emph{G}° = 68.4 kJ mol\textsuperscript{−1}

A Gibbs free energy of formation for liquid ethylene glycol is not given directly. However, the standard enthalpy of formation and the standard entropy are both given: C\textsubscript{2}H\textsubscript{6}O\textsubscript{2} (l) ∆\textsubscript{f}\emph{H}° = −460 kJ mol\textsuperscript{−1}, \emph{S}° = 163.2 J mol\textsuperscript{−1} K\textsuperscript{−1}

The formation reaction of ethylene glycol is as follows:

\begin{center}
2C (s) + 3 H\textsubscript{2} (g) + O\textsubscript{2} (g) → C\textsubscript{2}H\textsubscript{6}O\textsubscript{2} (l)
\end{center}

The standard entropies of the elements are well known: O\textsubscript{2} (g) \emph{S}° = 205.2 J mol\textsuperscript{−1} K\textsuperscript{−1}; H\textsubscript{2} (g) \emph{S}° = 130.7 J mol\textsuperscript{−1} K\textsuperscript{−1}; C (s) \emph{S}° = 5.7 J mol\textsuperscript{−1} K\textsuperscript{−1}.

Therefore, the standard entropy of formation for liquid ethylene glycol is: ∆\textsubscript{f}\emph{S}° = 163.2 J mol\textsuperscript{−1} K\textsuperscript{−1} − 205.2 J mol\textsuperscript{−1} K\textsuperscript{−1} − 3 × 130.7 J mol\textsuperscript{−1} K\textsuperscript{−1} − 2 × 5.7 J mol\textsuperscript{−1} K\textsuperscript{−1} = −445.5 J mol\textsuperscript{−1} K\textsuperscript{−1}

The standard Gibbs free energy of formation of liquid ethylene glycol is then:

\begin{center}
C\textsubscript{2}H\textsubscript{6}O\textsubscript{2} (l) ∆\textsubscript{f}\emph{G}° = ∆\textsubscript{f}\emph{H}° − \emph{T}∆\textsubscript{f}\emph{S}° = −460 kJ mol\textsuperscript{−1} − 298.15 K × (−445.5 J mol\textsuperscript{−1} K\textsuperscript{−1}) = −327.2 kJ mol\textsuperscript{−1}
\end{center}

For liquid water: H\textsubscript{2}O (l) ∆\textsubscript{f}\emph{G}° = −237.1 kJ mol\textsuperscript{−1}

With these standard Gibbs free energies of formation, the standard Gibbs energy change of the following reactions can be obtained:

\begin{center}
C\textsubscript{2}H\textsubscript{4} (g) + 2H\textsubscript{2}O (l) → C\textsubscript{2}H\textsubscript{6}O\textsubscript{2} (l) + H\textsubscript{2} (g)\\[2pt]
∆\emph{G}° = −327.2 kJ/mol − 2 × (−237.1 kJ mol\textsuperscript{−1}) − 68.4 kJ mol\textsuperscript{−1} = 78.6 kJ mol\textsuperscript{−1}
\end{center}

This process can be combined with the following two standard reduction potentials:

\begin{center}
2 H\textsuperscript{+} (aq) + 2 e\textsuperscript{−} → H\textsubscript{2} (g) \quad (\emph{E}° = 0.00000 V)\\[2pt]
OsO\textsubscript{4} (aq) + 2 H\textsubscript{2}O (l) + 2 e\textsuperscript{−} → OsO\textsubscript{2}(OH)\textsubscript{4}\textsuperscript{2−} (aq) \quad (\emph{E}° = 0.46 V)
\end{center}

The final process and its change in standard Gibbs energy is:

\begin{center}
OsO\textsubscript{4} (aq) + C\textsubscript{2}H\textsubscript{4} (g) + 4H\textsubscript{2}O (l) → OsO\textsubscript{2}(OH)\textsubscript{4}\textsuperscript{2−} (aq) + C\textsubscript{2}H\textsubscript{6}O\textsubscript{2} (l) + 2H\textsuperscript{+} (aq)\\[2pt]
∆\emph{G}° = 78.6 kJ mol\textsuperscript{−1} + 2 × 96485 C mol\textsuperscript{−1} × 0.000 V − 2 × 96485 C mol\textsuperscript{−1} × 0.46 V = −10.2 kJ mol\textsuperscript{−1}
\end{center}

The Henry constant of ethylene in water at 298.15 K is 21.7 kPa m\textsuperscript{3} mol\textsuperscript{−1}. Converted to the standard states (1 bar as the standard pressure, 1 mol kg\textsuperscript{−1} as the standard concentration, the standard density of water at 298.15 K as 997 kg/m\textsuperscript{3}), this is 21.6 bar kg mol\textsuperscript{−1}. Therefore, the equilibrium constant corresponding to the phase change of ethylene between aqueous and gaseous phases is 21.6 (dimensionless!), from which the standard Gibbs energy change of the process can be calculated:
\decompnote{Arithmetic check by the decompiler (not in the source): 21.7 kPa m\textsuperscript{3} mol\textsuperscript{−1} × 997 kg m\textsuperscript{−3} = 21 635 kPa kg mol\textsuperscript{−1} = 216 bar kg mol\textsuperscript{−1}, not 21.6; with 216 the next two results would be about −13.3 and −23.5 kJ mol\textsuperscript{−1}. Kept as printed.}

\begin{center}
C\textsubscript{2}H\textsubscript{4} (aq) → C\textsubscript{2}H\textsubscript{4} (g) \quad ∆\emph{G}° = −8.314 J mol\textsuperscript{−1} K\textsuperscript{−1} × 298.15 K × ln 21.6 = −7.6 kJ mol\textsuperscript{−1}
\end{center}

Combining this process with the one where gaseous ethylene is oxidized, one can obtain the standard Gibbs free energy of a process in which all the reactants are in the dissolved state:

\begin{center}
OsO\textsubscript{4} (aq) + C\textsubscript{2}H\textsubscript{4} (aq) + 4H\textsubscript{2}O (l) → OsO\textsubscript{2}(OH)\textsubscript{4}\textsuperscript{2−} (aq) + C\textsubscript{2}H\textsubscript{6}O\textsubscript{2} (l) + 2H\textsuperscript{+} (aq)\\[2pt]
∆\emph{G}° = −10.2 kJ mol\textsuperscript{−1} −7.6 kJ mol\textsuperscript{−1} = −17.8 kJ mol\textsuperscript{−1}
\end{center}

In the quantum chemical calculations, if the first reaction in the mechanisms is assumed to be an adduct formation between OsO\textsubscript{4} and C\textsubscript{2}H\textsubscript{4}, the product of which is the osmate ester OsO\textsubscript{4}⋅C\textsubscript{2}H\textsubscript{4}. The theoretical energy change for this process is reported as –21.2 kcal mol\textsuperscript{–1} (= –88.7 kJ mol\textsuperscript{–1})\textsuperscript{[46]} and –18.1 kcal mol\textsuperscript{–1} (= –75.7 kJ mol\textsuperscript{–1}).\textsuperscript{[47]} The same reaction in water would be:
\decompnote{The .docx has a soft hyphen (U+00AD) directly before ``–18.1''; it is invisible in print and left out here.}

\begin{center}
OsO\textsubscript{4} (aq) + C\textsubscript{2}H\textsubscript{4} (aq) → OsO\textsubscript{4}⋅C\textsubscript{2}H\textsubscript{4} (aq) \quad ∆\emph{G}°(1)
\end{center}

Then this intermediate reacts to give the final products in the following overall process:

\begin{center}
OsO\textsubscript{4}⋅C\textsubscript{2}H\textsubscript{4} (aq) + 4H\textsubscript{2}O (l) → OsO\textsubscript{2}(OH)\textsubscript{4}\textsuperscript{2−} (aq) + C\textsubscript{2}H\textsubscript{6}O\textsubscript{2} (l) + 2H\textsuperscript{+} (aq) \quad ∆\emph{G}°(2)
\end{center}

Thermodynamics requires that:

\begin{center}
∆\emph{G}°(1) + ∆\emph{G}°(2) = −17.8 kJ mol\textsuperscript{−1}
\end{center}

Undisputed observations show that OsO\textsubscript{4}⋅C\textsubscript{2}H\textsubscript{4} is an intermediate and not a final product. This means that ∆\emph{G}°(2) must be negative (otherwise, the process would not occur). Combining this inequity with the previous equation gives:

\begin{center}
∆\emph{G}°(1) > −17.8 kJ mol\textsuperscript{−1}
\end{center}

This conclusion is based on experimental data and observations only. It is definitive evidence that there can be no highly exergonic intermediate in the aqueous reaction.

One might argue that the basic medium in our experiments contributes significantly to the driving force of the reaction as the H\textsuperscript{+} (aq) produced is consumed (which is true), and this might allow an intermediate of much lower energy in the process. However, the adduct formation itself is pH-independent (no hydrogen ion, hydroxide ion or even water appears in it). Such a scenario would require the osmate ester as a stable final product in acidic solutions. This is certainly not the case, observation showed no detectable reaction at all below neutral pH.

\end{document}
